\section{Target-seeded inductive ideal completion}
\label{sec:ideals}

\subsection{The missing case after linear-space closure}

Consider $x'=x^2$, initial state $x=0$, and target $x=0$. The degree-one
stable-space algorithm correctly fails to keep $\Span\{x\}$: its pullback is
$x^2$, outside that vector space. But the ideal $\langle x\rangle$ is stable,
because $x^2=x\cdot x$. A polynomial multiplier is enough to prove the target.

This suggests a second lane. Instead of guessing a complete invariant basis and
its multipliers simultaneously, start with the desired target and close the
\emph{ideal it generates} under pullback. The auxiliary equations are then
forced by failed preservation checks. This is a target-directed specialization
of established algebraic invariant-set methods, not a new claim about the
foundations of invariant generation. In particular, the recent treatment by
Bayarmagnai, Mohammadi, and Pr\'ebet studies invariant-set computation and
polynomial-invariant testing with stronger geometric stopping criteria
\cite{bayarmagnai}. Here ordinary ideal membership is chosen because it returns
especially simple proof objects for Lean.

\subsection{A completion algorithm}

For the same unguarded rational-polynomial model as
Section~\ref{sec:spaces}, let $q$ be the target polynomial. Form the ascending
sequence of ideals
\begin{equation}\label{eq:idealchain}
 J_0=\langle q\rangle,\qquad
 J_{j+1}=J_j+\sum_{a=1}^{s}\langle F_a^*J_j\rangle.
\end{equation}
The notation on the right denotes the ideal generated by all pullbacks of
members of $J_j$. If $G_j$ generates $J_j$, it suffices to use the pullbacks of
those finitely many generators, because pullback is a ring homomorphism:
\[
 \left(\sum_i h_i g_i\right)\circ F_a
 =\sum_i (h_i\circ F_a)(g_i\circ F_a).
\]

The implementation adds one nonmember pullback at a time rather than computing
all of \eqref{eq:idealchain} eagerly. It keeps generators in their original
unreduced form, together with the transition word that produced them from $q$.
A Gr\"obner basis is an internal membership index, not the authoritative list of
invariant equations.

\begin{lstlisting}
G := [(target, emptyWord)]
repeat:
  for (g, word) in G:
    if g o initialization is not the zero polynomial:
      parameters := exactNonzeroEvaluation(g o initialization)
      return OrbitCounterexample(parameters, word)
  B, changeMatrix := tracedGroebnerBasis(polynomials(G))
  for transition F_a and (g, word) in G:
    p := g o F_a
    if p is not in ideal(G):
      append (p, prepend(a, word)) to G
      restart
  return initial-zero, pullback-membership, and target certificates
\end{lstlisting}

Prepending the new transition label is intentional. If
$g(x)=q(F_{a_m}(\cdots F_{a_1}(x)))$, then $g\circ F_b$ corresponds to executing
$b$ first and then $a_1,\ldots,a_m$. A noncommuting-transition regression test
checks this orientation independently by executing the returned word.

\begin{theorem}[Target completion for the unguarded model]\label{thm:ideal}
With exact ideal membership and without resource cutoffs, the algorithm
terminates. It either proves that $q$ vanishes on every reachable state from
$\phi(\Q^k)$, or returns a rational parameter tuple and a finite transition
word on which $q$ is nonzero.
\end{theorem}
\begin{proof}
Every added polynomial lies outside the current ideal, so each addition strictly
increases that ideal. The polynomial ring $\Q[x_1,\ldots,x_n]$ is Noetherian;
there can be only finitely many strict increases. Exact ideal membership is
effective by Gr\"obner-basis reduction. Thus either an initialization failure is
found or a pullback-stable ideal is reached.

At stable completion, all generators vanish initially, each generator's
pullback belongs to their ideal, and $q$ belongs to that ideal. Evaluating a
polynomial combination at a common zero produces zero. Induction on a transition
word therefore proves the target.

For the other outcome, every retained generator is an actual pullback of $q$
along its recorded word, not just a polynomial inferred by an untracked basis
change. If $g\circ\phi$ is nonzero, there exists a rational parameter tuple where
it evaluates nonzero. Concretely, if its degree in each parameter is at most
$D$, some point in $\{0,\ldots,D\}^k$ is nonzero: induction on the number of
variables reduces this fact to the bound on roots of a nonzero univariate
polynomial. Executing the recorded word at that parameterized initial state
therefore yields a counterexample to the target.
\end{proof}

This is a decision result for a \emph{specified polynomial equality} under
unconditional polynomial updates and unrestricted rational initialization
parameters. It is not a procedure for computing the full strongest invariant
ideal of an arbitrary program. Those are different problems; the known
hardness results for strongest polynomial invariants do not justify conflating
them \cite{mullner}.

\subsection{Why completion avoids the bilinear synthesis trap}

The two unknowns in a general preservation equation are the invariant family
and its polynomial multipliers. Completion never chooses both freely in one
constraint solve. At a membership check, the current generators are fixed;
the problem is to express a particular pullback in their ideal. If that fails,
the pullback itself becomes a new generator. All added equations are dictated
by the target and transition maps.

This has a useful proof-theoretic interpretation. A failed induction step asks
for precisely the missing equation needed to make that step close. Completion
accumulates those demands until they are mutually inductive. It is a specialized,
fully algebraic form of auxiliary-lemma discovery, rather than a language-model
proposal of unrelated lemmas.

For the coupled example, starting from $q=z-2y$ under $F_1$ produces
\[
 q\circ F_1=z-3y+x^2.
\]
This new equation is not an ideal multiple of $q$ and is added. The two equations
span the same linear family as $e_1,e_2$ from the earlier example, and both
transitions subsequently close. The implementation needs two generators and
one successful enlargement. Unlike the vector-space lane, it did not begin by
constructing every degree-two polynomial that vanishes initially.

\subsection{A genuinely polynomial action}

Take initialization $(x,y)=(t^2,t^3)$ and transition $(x,y)\mapsto(x^2,y^2)$.
The target $g=y^2-x^3$ satisfies
\[
 g\circ F=y^4-x^6=(y^2+x^3)g.
\]
The target has degree three, but its pullback has degree six. A degree-three
vector-space closure cannot retain $g$ merely by ignoring those higher-degree
terms. The ideal lane needs just one generator and the explicit multiplier
$y^2+x^3$.

More generally, $u^a-v^a=(u-v)\sum_{j=0}^{a-1}u^{a-1-j}v^j$ gives a family of
nonconstant action certificates for power-curve invariants. Twelve such fixtures
are included. They are constructed examples with known invariants, useful for
checking polynomial multiplication and substitution; they are not a held-out
benchmark of discovery performance.

\subsection{Traced membership, not a trusted CAS answer}

The search-side Gr\"obner implementation carries a representation
\[
 b_j=\sum_i T_{ji}g_i
\]
for every intermediate basis element. Forming an $S$-polynomial, reducing it,
and making it monic applies the same polynomial operations to its representation
row. If division gives $p=\sum_j u_j b_j$ with zero remainder, the desired
original-generator coefficients are $h_i=\sum_j u_jT_{ji}$.

Only the final identities are sent to the checker:
\begin{equation}\label{eq:idealcert}
 g_i\circ\phi=0,\qquad
 g_i\circ F_a=\sum_j h_{aij}(x)g_j(x),\qquad
 q=\sum_i c_i(x)g_i(x).
\end{equation}
The checker does not need to verify a Gr\"obner basis, rerun Buchberger, or trust
a CAS's zero-remainder flag. Conversely, a zero remainder relative to a new
basis is insufficient unless the change-of-generators expressions are retained.

For a production proposer, Singular's \code{lift} and \code{liftstd} facilities
are concrete candidates for obtaining these representations; their output must
still be converted to the original generator order and checked
\cite{singular}. SymPy is a convenient differential oracle for small cases, but
its ordinary basis reduction interface should not be mistaken for an already
integrated source-to-Lean certificate adapter \cite{sympy}.

\subsection{Counterexamples and conservative refusal}

The negative certificate is deliberately smaller than the positive one: it
contains initialization parameters and a finite word of transition indices.
The independent checker evaluates initialization, executes the transitions in
order using exact rationals, and verifies that the externally supplied target
is nonzero. It never trusts the intermediate polynomial that led the proposer
to that word.

Three fixtures produce such certificates: one fails after one step, another
after two steps, and a third requires the order $[1,0]$ of two noncommuting
maps. Reversing that last word is rejected. An ideal-generator cutoff separately
produces \code{unknown}, not a counterexample or an unprovability judgment.

When guards have been erased, an orbit counterexample is only a counterexample
to the enlarged model. To refute a source-level theorem, the adapter must replay
the same trajectory through the source guards and initialization domain. A
negative guard check invalidates the source-level counterexample, even though
the unguarded encoded observation remains correct.

\subsection{Where this lane belongs in the scheduler}

Noetherian termination is not a practical complexity guarantee. Gr\"obner bases,
coefficient sizes, and pullback degrees may become prohibitive. The prototype
limits generator count, pullback degree, division steps, and basis pairs; a
production worker additionally needs process-level memory and time limits,
because an expansion can become expensive before its completed degree is
inspected.

Run the linear-space lane first when $\binom{n+d}{d}$ is small and several
related targets are expected. Its result is reusable across all targets in the
computed space. Run ideal completion when a specific failed equality is available,
especially after the space lane reports degree escape. The latter is a useful
\emph{routing signal}, not a reason to truncate the polynomial representation.

The ideal lane is also a bridge to the cyclic compiler: a discovered family of
equations can serve as the simultaneous induction motive for a recursive
arithmetic helper. The two backends then cooperate through explicit local proof
obligations rather than through an unstructured bag of candidate equations.
